\subsection{Increasing homomorphisms of ordered groups}

Let G and \(G'\) be two ordered groups; among the homomorphisms f from the underlying additive group of G into that of \(G'\) , it makes sense to consider the increasing maps, that is those such that \(x \leqslant y\) implies \(f(x) \leqslant f(y)\) . Because of the relation \(f(y - x) = f(y) - f(x)\) it follows that the increasing homomorphisms from G into \(G'\) are characterised by the fact that the image of a positive element of G under such a homomorphism is a positive element of \(G'\) ; if P (resp. \(P'\) ) denotes the set of positive elements of G (resp. G'), this can be written \(f(P) \subset P'\) . It is clear that the canonical injection of a subgroup G into an ordered group G', and the projection of a product of ordered groups onto one of its factors, are increasing homomorphisms.

An isomorphism (VI, p.~2) f from an ordered group G onto an ordered group \(G'\) is a bijective homomorphism from G onto \(G'\) such that both f and the inverse homomorphism are increasing, in other words \(f(P) = P'\) .

It can happen that an isomorphism from the underlying group of G onto that of \(G'\) can be increasing without the inverse isomorphism also being increasing. This is the case, for example, if \(G = G'\) , if f is the identity map on G, and if \(P \subset P'\) but \(P \neq P'\) . Thus in Z we can take \(P'\) to be the set of (ordinary) positive integers and P to be the set of even positive integers.

(DIV) Let K be the field of rational functions \(\mathbf{F}_{2}(\mathbf{X})\) over the field \(F_{2}\) of order 2. The divisibility relations relative to the rings \(F_{2}[X] = A'\) and \(F_{2}[X^{2}, X^{3}] = A\) define two distinct ordered group structures on \(K^{*}\) , such that \(A \subset A'\) (they are ordered group structures since 1 is the only unit in A and the only unit in \(A'\) ).

